\appendix
\section{Appendix}
\label{sec:appendix}

\subsection{Estimator validation against planted ground truth}
\label{sec:app-planted}

Three of the measurements in this thesis rest on estimators built for the purpose rather than taken
off the shelf. Each was validated by planting a known answer and checking recovery --- in both
directions, since an estimator that only ever returns the answer one hopes for is not evidence.

\begin{table}[h]
\centering
\begin{tabular}{llrr}
\toprule
Estimator & Planted truth & Recovered & \\
\midrule
transfer coefficient $\Tcoef$ & $\Tcoef = 1.00$ (no interaction) & 0.999 & \\
 & $\Tcoef = 0.70$ & 0.703 & \\
 & $\Tcoef = 0.40$ & 0.404 & \\
\midrule
channel gate & a channel that genuinely predicts & $+0.397$ (sd 0.029) & over 24 worlds \\
 & a channel carrying nothing & $-0.016$ (sd 0.091) & over 24 worlds \\
\midrule
$\kappa$-structure & a consistent per-line modulation & excess $+1.085$ & \\
 & an idiosyncratic interaction & excess $+0.001$ & \\
\bottomrule
\end{tabular}
\caption{All three recover planted answers in both directions.}
\label{tab:planted}
\end{table}

Two of these validations changed the estimators rather than merely endorsing them. The channel
gate's count-matched null originally drew its random partners from a pool \emph{excluding} the
channel's own selections, which depletes the null of exactly the partners the channel chose and
anti-correlates the two arms. And the useless-channel gap has a standard deviation of roughly $0.09$
on a single small world, so the original validation --- asserting that one draw fell within $0.10$ ---
tested luck rather than correctness; it now averages 24 independent worlds. That variance is also why
the reported channel results are gated on a clustered interval rather than on a point estimate.

\subsection{Absent-gene convention}
\label{sec:app-convention}

A predicted cell sentence lists only expressed genes, so some panel genes are unranked and a
convention is required. Two were carried throughout: \texttt{worst} assigns absent genes the last
rank $P$, \texttt{francesca} a fixed middle rank $P/2$. The choice does not matter --- spike-in
$\DEdr$ reads $0.952$ under the first and $0.954$ under the second, and no ceiling or baseline in
this thesis moves by more than $0.01$ between them. All figures in the body are quoted under
\texttt{worst}.

\subsection{Swap-distance sweep}
\label{sec:app-sweep}

The stratified scramble of Section~\ref{sec:methods-controls} was preceded by a continuous version:
$\gap$ measured as a function of the response distance between the real and the swapped-in drug,
binned into quintiles.

\begin{table}[h]
\centering
\begin{tabular}{lrr}
\toprule
Bin (swap distance) & single-cell model & optimal-transport model \\
\midrule
0 --- nearest (2.75--4.03) & $-0.046$ & $-0.001$ \\
1 (4.03--4.60) & $+0.060$ & $+0.031$ \\
2 (4.60--4.93) & $-0.006$ & $+0.017$ \\
3 (4.93--5.32) & $-0.004$ & $+0.043$ \\
4 --- farthest (5.32--6.25) & $-0.019$ & $-0.005$ \\
\midrule
trend, corr(distance, $\gap$) & $+0.051$ & $+0.012$ \\
\bottomrule
\end{tabular}
\caption{No bin in either model has an interval excluding zero, and the trend correlations are
$\approx 0$. 200 (drug, swap) pairs over 20 cell lines, $n=40$ per bin.}
\label{tab:sweep}
\end{table}

Two features make this an argument rather than five nulls. The single borderline value sits in a
\emph{middle} bin, whereas a genuine dissimilarity effect would be monotone and peak in the far bin.
And the comparison is paired: the model's own $\NIR$ drifts across bins from $0.497$ to $0.326$,
because drugs that are far from everything populate the far bins, but $\gap$ uses the same drug on
both arms so that drift cancels.

The intervals are wide --- $n=40$ per bin gives half-widths of $0.06$--$0.10$ --- so this sweep
excludes a large effect rather than establishing a null, which is why the stratified scramble
replaced it in the body.

\subsection{Training configuration}
\label{sec:app-training}

Every arm was cold-started from the same base checkpoint with identical hyperparameters, so that the
target or objective is the only variable. One exception is material and is stated wherever the arm
is quoted.

\begin{table}[h]
\centering
\begin{tabular}{llll}
\toprule
Arm & Target & Training & Note \\
\midrule
single cell & the observed treated cell & 1 epoch & \\
consensus & group pseudobulk & \textbf{$\approx 60\%$ of an epoch} & tier-1 $\NIR$ unscoreable \\
optimal transport (T2) & barycentric, $\varepsilon = 0.5$ & 1 epoch, 23{,}842 steps & \\
residual & signed residual signature & 1 epoch, 15{,}340 steps & \\
residual (holdout) & as above, 738 conditions withheld & 1 epoch & \\
\bottomrule
\end{tabular}
\caption{Learning rate $10^{-5}$, batch size 1, gradient accumulation 16, maximum length 8192, bf16,
weight decay 0.01, warmup ratio 0.03 throughout. The consensus arm stopped early on a flat loss;
since it is one endpoint of the $\varepsilon$-ladder, the ladder's negative rests on the other two.}
\label{tab:trainconfig}
\end{table}

\subsection{Calibration self-test}
\label{sec:app-selftest}

The calibration framework of Section~\ref{sec:methods-metrics} is itself checked, since a metric
audit whose instrument is untested proves nothing. Beyond deterministic unit checks --- a perfect
prediction must reach the metric's maximum, a zero shift and a flat truth must return no score ---
the self-test asserts the exploit under a \emph{second, independent} zero-information baseline: a
leave-one-out mean scores $\DEdr = 1.000$ while scoring $\NIR = 0.000$. Two unrelated constructions
both saturating the standard metric, and both scoring zero on the discrimination metric, establishes
the exploit as a property of the metric rather than of any one baseline.
